\section{Initial increments at every Sobolev order}
\label{sec:initial-summability}

We use the joined displacement-history and forward inverse of
Lemma~\ref{lem:packet-joined-high}. When the activation time is positive,
the forced high coefficients generally have nonzero initial velocities;
all of them enter the estimates below.

Fix the dimension $d$ and a legal integer base Sobolev order $s$ for the
packet estimates. Constants below may depend on $d$ and on a fixed
derivative order, but the frequency condition will not depend on that
derivative order. Write $\mathcal P\ge1$ for a bound on the actual source primitives,
including the coefficient and inverse costs, the terminal datum, the
reciprocal phase scale, the fixed mean cutoff cost, and the common Gevrey
radius. Section~\ref{sec:parameter-envelope} bounds these primitives for the
constructed sources, uniformly along the prescribed scale sequence.

\subsection{The actual three sums at time zero}
Put $\kappa=k^{-1}$ and let $N$ be the chosen truncation. Since the exact
nonlinear correction has zero initial value and the parent flow is initially
the identity, the joined finite construction gives
\begin{align}
 u_{\mathrm{osc},0}(a)
 &=\rho\sum_{p=1}^{N}
    \bigl[\kappa^p A_p+\kappa^{p+1}C_p\bigr]
       (0,a/\rho,km_0\cdot a/\rho),\label{eq:initial-osc-sum}\\
 u_{\mathrm{mean},0}(a)
 &=\rho\sum_{p=2}^{N}\kappa^p B_p(0,a/\rho),
 \qquad
 \Delta u_0=u_{\mathrm{osc},0}+u_{\mathrm{mean},0}.
 \label{eq:initial-mean-sum}
\end{align}
Here $C_p$ is the full ambient antisymmetric-potential corrector and $B_p$
is independent of the angle. The estimates include every high coefficient, its tensor corrector,
and every mean coefficient.
Formula~\eqref{eq:initial-osc-sum} extends
\cite[\S3.9, (3.60)]{OAI} to the full ambient tensor construction.

Let $\alpha>0$ be the selected primary amplitude. The joined profile is
\[
 \gamma(t)=\alpha\quad(0\le t\le t_0),\qquad
 \gamma(t)=\alpha g(t)\quad(t_0\le t\le T),\qquad g(t_0)=1,
\]
and put $H_0=\max(1,\alpha,\sup\alpha g)$. The actual coefficient induction
assigns the oscillatory profiles $\gamma H_0^{2p-2}$ and the mean profiles
$H_0^{2p-2}$.

The factor $\alpha$ on the history interval can be checked directly.
The primary terminal displacement is $\alpha$ times the normalized compact
datum, and the displacement inverse is linear. In every higher zero-mean
force, terms containing only means cancel under angle-mean subtraction.
Every remaining term contains a previous oscillatory factor or a derivative
of one. The term involving the newly solved mean is
$-(m\cdot B_p)\partial_\theta A_1$, which also contains this factor.
Extra oscillatory factors are bounded by the indicated powers of $H_0$.
The actual history inverse and the angular tensor primitive therefore
retain $\alpha$ at time zero. This factor is preserved at each step of the graded construction
by linearity of its high solver.

Suppose first that $t_0>0$. In fixed transverse coordinates, every
forced history solves
$\xi(0)=\xi(t_0)=0$ and has velocity $A=FR_\perp\xi_t$.
The primary instead has $\xi(0)=0$ and
$\xi(t_0)=\alpha\chi_{\rm in}(y)f_\delta(\theta)\xi_T$, where $\xi_T$ is the
actual selected terminal coordinate. It is obtained from the physical
central endpoint by the actual inverse frame; its norm is included in $\mathcal P$.
Subtracting the affine extension
$(t/t_0)\alpha\chi_{\rm in}f_\delta\xi_T$ reduces this primary problem to the
same zero-endpoint inverse. All spatial and phase words of the forced
problem still have zero endpoint values. For the primary, differentiating
this explicit affine extension supplies exactly the differentiated
boundary datum; the differentiated boundary traces are therefore explicit.
The forward solution starts with the actual trace $FR_\perp\xi_t(t_0)$.
It agrees there with the history solution, and the common velocity equation
then gives equality of the time derivatives at the junction. These are
the boundary conditions of Lemma~\ref{lem:packet-joined-high}.
When $t_0=0$, the primary instead has its prescribed forward datum
and the forced high fields have zero initial velocity. The forward
estimates give the same bounds below, with no history inverse or
reciprocal factor $t_0^{-1}$. This applies both to the exceptional seed
and to the first normal stage $j=J$; positive history begins at $j=J+1$.

For completeness, the all-word estimate uses this same inverse. After
translation in the label and angle, write the equation as
$L(a)u(a)=f(a)$. At a fixed legal base Sobolev order let its inverse cost
be $M$, and let the known coefficient derivative costs be
$B R_c^j(j!)^2$. Differentiating a word of length $n$ gives the exact
Leibniz identity
\[
 L\partial^\omega u=\partial^\omega f-
 \sum_{\varnothing\ne S\subset\{1,\ldots,n\}}
 (\partial^{\omega_S}L)\partial^{\omega_{S^c}}u.
\]
The homogeneous boundary space is unchanged, so the same inverse applies
to each term. In the block norm this yields a binomial convolution.
For $b\ge0$,
\[
 \binom nj(j!)^2((n-j+b)!)^2\le((n+b)!)^2.
\]
Consequently a forcing bound $F R^{n+b}((n+b)!)^2$ gives
$2MF R^{n+b}((n+b)!)^2$ for the solution as soon as
$R\ge R_c(1+2MB)$: the positive-order convolution costs at most
$MB\sum_{j\ge1}(R_c/R)^j\le1/2$.
The initial finite base inverse is obtained by the usual finite Sobolev
commutator recursion. Its order is fixed by $d$, not by the packet grade.
Time traces follow from the actual $H^1$ equation and cost a polynomial
in the reciprocal history length and known frame bounds. The same estimate
on the forward interval uses its proved relative propagator and the actual
history trace. This proves preservation of the unknown forcing radius;
known coefficient conversion enlarges $R_c$ only once.

The displacement argument uses the whole transverse space of dimension
$d-1$. If $R_\perp$ has orthonormal columns, $\eta=FR_\perp\xi$, and
$\eta(t_0)=0$, then
\[
 \int_0^{t_0}\bigl(|\eta_t|^2-\langle H\eta,\eta\rangle\bigr)\,dt
 \ge \left(1-\frac{K_+t_0^2}{2}\right)
                       \int_0^{t_0}|\eta_t|^2\,dt
 \ge\frac12\int_0^{t_0}|\eta_t|^2\,dt.
\]
The one-sided Poincar\'e inequality here follows by integrating
$\eta(t)=-\int_t^{t_0}\eta_t(r)\,dr$. Conversion to $\xi$ uses only the
known bounds for $F^{-1}$ and its derivative. Thus the inverse cost is
polynomial in those bounds, and this step has no three-dimensional
cross-product dependence.

\subsection{Separating derivative order from expansion order}
For a fixed integer $m\ge0$, the coefficient estimates need derivatives
only through $m+s+1$. Their shift is at most $100p$; replacing this fixed
coefficient by a larger dimensional constant makes no change to the
argument. For $b=m+s+1$,
\[
 (b+100p)!\le2^{b+100p}b!(100p)!,\qquad
 ((b+100p)!)^2\le C_b(Cp)^{200p}.
\]
Also $R^{b+100p}=R^b(R^{100})^p$. Consequently all dependence on $m$
can be put in a prefactor $C_{d,m}\mathcal P^{c_{d,m}}$, while the
expansion-order bound has the form
\begin{equation}
 B_N^{p+1},\qquad B_N=\mathcal P^{c_d}(C_dN)^{L_d}.
 \label{eq:initial-grade-base}
\end{equation}
with one finite exponent $L_d$ independent of $m$.
For the displayed shifts one may take $L_d=300$ after enlarging constants.
Choose $N=\lfloor k^\vartheta\rfloor$ with
$L_d\vartheta<1/100$. A single fixed source-to-frequency guard then
ensures $B_N\le k^{1/100}$.

At time zero the cylinder Sobolev norms of $A_p$ and $C_p$ are bounded by
\[ \alpha C_{d,m}\mathcal P^{c_{d,m}}B_N^{p+1}. \]
The mean has the same bound without $\alpha$. For $k\ge4$,
\[
 \sum_{p=1}^{N}k^{-p}B_N^{p+1}\le C,
 \quad
 \sum_{p=1}^{N}k^{-p-1}B_N^{p+1}\le C,
 \quad
 \sum_{p=3}^{N}k^{-p}B_N^{p+1}\le Ck^{-2}.
\]
These follow by summing the geometric series with ratio $k^{-99/100}$.
In the last sum its first power is $k^{-3+4/100}\le k^{-2}$.
The grade-two mean is estimated separately at its fixed order, giving
$C_{d,m}\mathcal P^{c_{d,m}}k^{-2}$ as well. This fixed-order separation prevents
an artificial loss of a small power of $k$ in the mean estimate.

For $T_kf(y)=f(y,km_0\cdot y)$ and $|m_0|=1$, the chain rule and
$H^1(\mathbb T)\hookrightarrow L^\infty(\mathbb T)$ give
\[
 \|T_kf\|_{H^m(\R^d)}
       \le C_{d,m}(1+k)^m
                    \|f\|_{H^{m+1}(\R^d\times\mathbb T)}.
\]
Indeed each spatial derivative is a product of the commuting operators
$\partial_{y_i}+km_{0,i}\partial_\theta$. Apply the periodic one-dimensional
estimate to each resulting derivative at fixed $y$, then integrate in $y$.
Exactly one additional angle derivative is needed. For the mean there is
no graph loss because the field is angle independent.
Finally, physical scaling multiplies the $L^2$ norm of an order-$r$
derivative by $\rho^{1+d/2-r}\le\rho^{-m}$ for $r\le m$ and
$0<\rho\le1$. We obtain the actual estimates
\begin{align}
 \|u_{\mathrm{osc},0}\|_{H^m}
 &\le C_{d,m}\rho^{-m}k^m\alpha \mathcal P^{c_{d,m}},\label{eq:joined-initial-high}\\
 \|u_{\mathrm{mean},0}\|_{H^m}
 &\le C_{d,m}\rho^{-m}k^{-2}\mathcal P^{c_{d,m}}.
 \label{eq:joined-initial-mean}
\end{align}
The selected central amplitude satisfies
\[
 \alpha\le \mathcal P^{c_d}e^{-b x}
\]
for a fixed $b>0$, by the actual central target growth, with
the normalized terminal and reciprocal phase costs included in $\mathcal P$.
In the scalar indexing below, $n=j-J$ corresponds to normal stage $j$
of Section~\ref{sec:iteration}: $a_n=j$, $x_n=x_{j-1}$,
$k_n=k_j$, and $\rho_n=\rho_j$. The inherited geometric parameter
is therefore $x_n$, while the new target is $x_{n+1}$. The amplitude
factor is $e^{-b x_n}$. The required exponential gain is present for
all initial high coefficients, including those produced by positive
history intervals.

\subsection{One scale sequence gives all orders of summability}
\begin{lemma}[Scalar summability]\label{lem:initial-scalar-sum}
Fix $J\ge2$, $X\ge1$, and define
\[
 a_n=J+n,\quad x_0=X,\quad x_{n+1}=a_n^2x_n,
 \quad k_n=e^{x_n/a_n^2},\quad \rho_n=e^{-x_n/a_n^{7/2}}.
\]
Here $\rho_n$ is the physical packet length; it is distinct from the
fixed physical mean-cutoff parameter $\ell_*$ and from the Gevrey
weight radius used in the cylinder estimates. Let
\[
 \mathcal P_n=C a_n^p x_n^q
           \exp\left(\frac{c x_n}{(a_n-1)^3}\right)
\]
with fixed $C\ge1$, $c\ge0$, $p,q\in\N$ and $b>0$. For every fixed
$m,\nu\in\N$, both series
\[
 \sum_n\rho_n^{-m}\mathcal P_n^\nu k_n^m e^{-b x_n},
 \qquad
 \sum_n\rho_n^{-m}\mathcal P_n^\nu k_n^{-2}
\]
converge.
\end{lemma}
\begin{proof}
We have $x_n\ge X4^n$ and
$\log x_n\le\log X+2n\log(J+n)$. Therefore, for every fixed $r\ge0$,
\[
 \frac{a_n^r(1+\log a_n+\log x_n)}{x_n}\longrightarrow0.
\]
The logarithm of the first summand, divided by $x_n$, tends to $-b$:
its other exponential terms are $m/a_n^2$, $m/a_n^{7/2}$ and
$c\nu/(a_n-1)^3$, and its logarithmic terms vanish. Eventually the summand
is at most $e^{-b x_n/2}$, a summable sequence.

For the second summand, divide its logarithm by $x_n/a_n^2$. The result
tends to $-2$, since
\[
 ma_n^{-3/2}\to0,\qquad
 c\nu\frac{a_n^2}{(a_n-1)^3}\to0,
\]
and the logarithmic terms again vanish. Eventually it is at most
$e^{-x_n/a_n^2}\le e^{-n}$, proving convergence.
\end{proof}

The choices of $J,X$ in this lemma are independent of $m$. Only the
starting index of the convergence comparison depends on $m$; the
construction itself uses the same scales for every Sobolev order.
The same proof permits a localization scale whose reciprocal is bounded
by a polynomial times $\exp(c_0x_n/a_n^\beta)$ for any fixed $\beta>2$.
The quantitative localization estimate is therefore part of the
source construction.

To check that the prescribed frequency can satisfy all construction
guards, consider any finite list $A\mathcal P_n^\nu\le k_n^\theta$ with
$\theta>0$. Choose $J$ so that
$c\nu a^2/(a-1)^3\le\theta/2$ for all $a\ge J$ and every guard.
Writing $x_n=Xb_n$ with $b_n\ge J^{2n}$, the sequences
$a_n^2/b_n$, $a_n^2\log a_n/b_n$ and $a_n^2\log b_n/b_n$ are bounded.
Hence a single increase of $X$ pays all the remaining logarithmic terms,
uniformly in $n$. This proves the guards at the actual $k_n$.
Only finitely many construction guards enter at fixed $d$; the fixed-$m$
prefactors in \eqref{eq:joined-initial-high}--\eqref{eq:joined-initial-mean}
do not create new construction guards.

\subsection{The smooth compact limit}
Section~\ref{sec:parameter-envelope} bounds the stage-$n$ source
primitives by $\mathcal P_n$ and verifies the prescribed localization
scale for the joined construction. The preceding estimates consequently give
\[
 \sum_n\|\Delta u_n(0)\|_{H^m}<\infty
 \qquad\text{for every integer }m\ge0.
\]
For each spatial order $r$, take $m>r+d/2$. The legal-dimensional Sobolev
embedding then gives uniform convergence of the derivative series through
order $r$. Include the base solution and the one exceptional seed increment in
$u_{\mathrm{seed},0}$. Consequently the single sum
$u_0=u_{\mathrm{seed},0}+\sum_n\Delta u_n(0)$ is smooth, has the termwise
derivatives, and is divergence free. The fixed physical mean-cutoff
construction puts the supports of all these genuine initial increments
in one compact ball. Enlarging it to contain the base support proves
$u_0\in C^\infty_{c,\sigma}(\R^d)$.

All estimates concern the joined family of
Section~\ref{sec:parameter-envelope} at the prescribed frequencies.
